%!TEX root = main.tex

In this section, we describe the main ideas used in the formalization of HiFIRRTL.

Let us focus on a HiFIRRTL module, say $M$.

A structure called circuit diagram is used in the formalization to represent HiFIRRTL programs.

\begin{definition}
A \emph{circuit diagram} $D$ is a triple $(ss, ce, cm)$, where $ss$ is a sequence of statements, $ce$ is a component environment that maps each component to its type, and $cm$ is a connect mapping that maps each component to an expression that is connected to it. 
\end{definition}

A HiFIRRTL module $M$ will be first transformed into a circuit diagram $D_M = (ss, ce, cm)$, where $ss$ is the sequence of statements in $M$, $ce$ is the component environment obtained from the variable declarations in $M$, possibly with the types of some variables missing or partial, and $cm$ is the trivial mapping that maps each component to $\bot$ (Intuitively, this means that the connects of the components have not be determined yet). The missing or partial  types in $ce$ will be completed by the compilation passes during the transformation of $M$ into a LoFIRRTL module. Moreover, the connect mapping $cm$ will be updated by the ExpandWhens compilation pass. Furthermore, the statement sequence $ss$ will be updated by ExpandConnects, ExpandWhens, LowerTypes, and RemoveReset, while they keep unchanged during the other compilation passes.  

